\section{Implementation details of all teleoperation methods}
\label{app:methods}

All methods implement one interface: \texttt{reset(q0)} and
\texttt{solve(targets, dt)} mapping absolute EE pose targets (robot base
frame, both arms) to a 10-DoF joint command in a fixed joint order. The
\texttt{MethodIKAdapter} wraps any of them behind the production solver
interface and applies a joint-space rate limiter
\(\lvert\Delta q\rvert \le \dot q_{\max}\,\Delta t\) after every solve
(\(\dot q_{\max}=3\)~rad/s in simulation, 2~rad/s on hardware)~--- the
benchmark showed \texttt{scipy\_ls} and \texttt{mink} can jump between IK
branches near the workspace edge, and the limiter bounds the jump rate.
Grippers are excluded from IK everywhere (locked in the reduced models).

\subsection{\texttt{pink\_full} — production QP, full orientation cost}
\label{app:pink-full}

Differential IK with Pink over the reduced Pinocchio model. Tasks: a
6D \texttt{FrameTask} per EE (position cost 1.0, orientation cost 0.75,
task gain 0.4, Levenberg–Marquardt damping 0), a \texttt{DampingTask}
(cost 0.25), and a vectorized posture task (zero cost by default). The QP
solves for joint velocities under configuration and velocity limits
(solver \texttt{quadprog}, Tikhonov damping \(10^{-12}\)), then integrates
with \(\Delta t = 10\)~ms. This is the tracker the production thread has
always used; on the 5-DoF arm its full orientation cost trades position
error for unreachable-orientation error (Section~\ref{sec:exp-tracking}).
In production the armplane construction (Section~\ref{sec:armplane})
feeds it only reachable orientations, which largely neutralizes the
trade-off.

\subsection{\texttt{pink\_relaxed} — QP with relaxed orientation}
\label{app:pink-relaxed}

Identical QP with the orientation cost lowered to 0.05: orientation becomes
a soft preference and position dominates. Best all-round tracker of the
original benchmark and the current recommended production setting when
paired with the armplane mapping.

\subsection{\texttt{dls} — damped least squares}
\label{app:dls}

Per arm and per tick, the 6D task error
\(e = [\,p^\ast - p;\ w_o \log(R^\ast R^\top)\,]\) (orientation weight
\(w_o=0.5\)) and the site Jacobian \(J\) (weighted likewise) give
\(\Delta q = J^\top (J J^\top + \lambda^2 I_6)^{-1} e\) with
\(\lambda = 0.05\), applied with per-tick gain 0.4, velocity clamp
3~rad/s, and joint-range clamp. Runs directly on the MuJoCo model
(vr-teleop-kit style), so it also cross-checks the Pinocchio methods
against the simulator's kinematics.

\subsection{\texttt{mink} — MuJoCo-native QP}
\label{app:mink}

The mink library's differential IK over the MuJoCo model: per-EE 6D frame
tasks with the production cost ratio, solved as a QP with joint limits and
an added velocity cap matching \texttt{dls}. Behaviorally close to
\texttt{pink\_full} (same task structure, different modeling backend).

\subsection{\texttt{scipy\_ls} — bounded nonlinear least squares}
\label{app:scipy-ls}

Per tick, a \texttt{scipy.optimize.least\_squares} solve (Trust Region
Reflective) over the arm joints with joint bounds, minimizing the stacked
position and (weighted) orientation residuals from the previous solution as
the initial guess. Accurate but the most expensive per tick, and prone to
branch jumps near the workspace edge~--- the adapter's rate limiter is
mandatory.

\subsection{\texttt{telegrip} — split IK, direct wrist}
\label{app:telegrip}

Port of Telegrip's control scheme (Section~\ref{sec:telegrip}). Wrist
first: \(q_4\) from the affine elevation model
\(\theta = e_0 + s_2 q_2 + s_3 q_3 + s_4 q_4\) evaluated with the currently
commanded \(q_2, q_3\); \(q_5\) from the affine roll model
\(\psi = r_0 + s_5 q_5\) (previous roll held when the tip is within
\(5.7^\circ\) of vertical). Coefficients calibrated by finite differences
on the simulator's kinematics at construction; slopes are validated to be
exactly \(\pm1\) and the affine models to \(<1^\circ\) by unit tests.
Then position: DLS on the \(3\times3\) position Jacobian of
\(q_1..q_3\) with \(\lambda=0.05\), per-tick gain 0.6. All joints
velocity-clamped (3~rad/s) and range-clamped. Differences from upstream
Telegrip: PyBullet is replaced by the benchmark scene's kinematics; the
wrist angles are recovered from the absolute target rotation rather than
integrated from controller deltas (equivalent under the clutch mapping,
but stateless); and the shoulder pan participates in the position DLS
rather than being wrapped separately~--- the \(\pm110^\circ\) pan range
makes the differential step's implicit wrapping sufficient.

\subsection{Production input processing (common to all methods)}

One-Euro filtering of controller transforms
(\(f_{c,\min}=0.8\)~Hz, \(\beta=5.0\), \(f_{c,d}=0.9\)~Hz; position and
quaternion filtered separately with a sign-consistency guard); grip
threshold 0.9 on both hands for the clutch; per-session handle-axis
calibration at first grip; orientation blend-in over 1~s after activation;
optional height lock (freezes target \(z\) per arm) and mirror mode (swaps
and laterally reflects the hands); workspace-envelope policy
(Section~\ref{sec:envelope}) applied after the height lock, immediately
before targets reach the IK.
